% 量子力学1 量子力学1-13 問題1（089-step） 解いた日 2026-11-21
%% kotae
\toi{1}{○}{記述}
（図：$x<0$ で右向きの入射波 $\phi_{1}=\ee^{\ii kx}$ と左向きの反射波 $r\ee^{-\ii kx}$，$x>0$ で透過波 $\phi_{2}=t\ee^{\ii qx}$。段差は $x=0$）

$\phi_{1}$ について
\begin{align*}
&-\frac{\hbar^{2}}{2m}(\ii k)^{2}\ee^{\ii kx}=E\ee^{\ii kx}\\
&\therefore\ \frac{\hbar^{2}k^{2}}{2m}=E\qquad k=\frac{\sqrt{2mE}}{\hbar}
\end{align*}
$\phi_{2}$ について
\begin{align*}
&-\frac{\hbar^{2}}{2m}(\ii q)^{2}t\ee^{\ii qx}+V_{1}t\ee^{\ii qx}=Et\ee^{\ii qx}\\
&\therefore\ \frac{\hbar^{2}q^{2}}{2m}=E-V_{1}\\
&\qquad q=\frac{\sqrt{2m(E-V_{1})}}{\hbar}
\end{align*}
\begin{align*}
j&=\frac{\hbar}{2m\ii}\Bigl((\ee^{-\ii kx}+\ayamari{r}{r^{*}}\ee^{\ii kx})\\
&\qquad\times(\ii k\ee^{\ii kx}-\ii kr\ee^{-\ii kx})\\
&\qquad-(\ee^{\ii kx}+r\ee^{-\ii kx})\\
&\qquad\times(-\ii k\ee^{-\ii kx}+\ii k\ayamari{r}{r^{*}}\ee^{\ii kx})\Bigr)\\
&=\frac{\hbar k}{2m}\Bigl((\ee^{-\ii kx}+r\ee^{\ii kx})(\ee^{\ii kx}-r\ee^{-\ii kx})\\
&\qquad-(\ee^{\ii kx}+r\ee^{-\ii kx})(-\ee^{-\ii kx}+r\ee^{\ii kx})\Bigr)\\
&=\frac{\hbar k}{2m}\Bigl((1+r\ee^{2\ii kx}-r\ee^{-2\ii kx}-r^{2})\\
&\qquad-(-1+r\ee^{2\ii kx}-r\ee^{-2\ii kx}+r^{2})\Bigr)\\
&=\frac{\hbar k}{2m}\cdot2(1-r^{2})\quad\text{干渉項は打ち消し合う}\\
&=\frac{\hbar k}{m}(1-r^{2})\qquad(x<0)
\end{align*}
\naze{$\phi^{*}$ では係数も複素共役にするので $\phi^{*}=\ee^{-\ii kx}+r^{*}\ee^{\ii kx}$ である（$x>0$ も $t^{*}\ee^{-\ii qx}$）。以降はこの誤りを引き継ぎ，正しくは $j=\hbar k(1-\abs{r}^{2})/m$（$x<0$），$j=\hbar q\abs{t}^{2}/m$（$x>0$）である。干渉の項は $r^{*}\ee^{2\ii kx}-r\ee^{-2\ii kx}$ で純虚数なので，$\ii k$ を掛けると実数になり虚部に残らない。}
\begin{align*}
j&=\frac{\hbar}{2m\ii}\bigl(\ayamari{t}{t^{*}}\ee^{-\ii qx}\cdot\ii qt\ee^{\ii qx}\\
&\qquad-t\ee^{\ii qx}\cdot(-\ii q)\ayamari{t}{t^{*}}\ee^{-\ii qx}\bigr)\\
&=\frac{\hbar q}{2m}(t^{2}+t^{2})\\
&=\frac{\hbar q}{m}t^{2}
\end{align*}
\toi{2}{◎}{}
\begin{align*}
&\lim_{x\to-0}\phi(x)=1+r\\
&\phi(0)=t\qquad\therefore\ 1+r=t\quad\cdots\text{①}\\
&\phi'(x)=\ii k\ee^{\ii kx}-\ii kr\ee^{-\ii kx}\quad(x<0)\\
&\phi'(x)=\ii qt\ee^{\ii qx}\quad(x>0)\\
&\lim_{x\to-0}\phi'(x)=\ii k(1-r)\\
&\phi'(0)=\ii qt\qquad k(1-r)=qt\quad\cdots\text{②}
\end{align*}
①，②より
\begin{align*}
&k+kr=kt\\
&k-kr=qt\\
&2k=(k+q)t\\
&t=\frac{2k}{k+q}
\end{align*}
\begin{align*}
r&=\frac{k-q}{2k}t\\
&=\frac{k-q}{2k}\cdot\frac{2k}{k+q}=\frac{k-q}{k+q}
\end{align*}
\toi{3}{◎}{}
\begin{align*}
R&=\frac{\abs{j_{\mathrm{ref}}}}{j_{\mathrm{in}}}=\frac{\abs{-\frac{\hbar k}{m}r^{2}}}{\frac{\hbar k}{m}}=r^{2}\\
&=\left(\frac{k-q}{k+q}\right)^{2}
\end{align*}
\begin{align*}
T&=\frac{j_{\mathrm{tr}}}{j_{\mathrm{in}}}=\frac{\frac{\hbar q}{m}t^{2}}{\frac{\hbar k}{m}}=\frac{q}{k}t^{2}\\
&=\frac{4kq}{(k+q)^{2}}
\end{align*}
\begin{align*}
R+T&=\frac{k^{2}-2kq+q^{2}+4kq}{(k+q)^{2}}\\
&=\frac{(k+q)^{2}}{(k+q)^{2}}=1
\end{align*}
$T\neq\abs{t}^{2}$ は $x\ge0$ と $x<0$ で波数が違うから。
\toi{4}{△}{方針}
\begin{align*}
R&=\frac{k^{2}+q^{2}-2kq}{k^{2}+q^{2}+2kq}\\
&=\frac{2mE+2m(E-V_{1})-4m\sqrt{E(E-V_{1})}}{2mE+2m(E-V_{1})+4m\sqrt{E(E-V_{1})}}\\
&=\frac{2E-V_{1}-2\sqrt{E(E-V_{1})}}{2E-V_{1}+2\sqrt{E(E-V_{1})}}\\
&=\frac{6-2\sqrt{5}}{6+2\sqrt{5}}\\
&=\frac{(6-2\sqrt{5})^{2}}{16}=\ayamari{0.145}{0.146}
\end{align*}
\naze{$(6-2\sqrt{5})^{2}/16=\num{0.14590}$ なので，3桁では次の桁を四捨五入して $\num{0.146}$ である（切り捨てたので $R+T=\num{0.999}$ になる）。}
\[ T=1-R=0.854 \]
古典粒子では $\ayamari{R=1,\ T=0}{R=0,\ T=1}$
\naze{$E>V_{1}$ の古典の粒子は段差で減速するだけで必ず越えるので，$R=0$，$T=1$ である。量子力学では越えられるエネルギーでも $\num{0.146}$ の割合で反射する，というのが比べる要点である。}
%% kizuki
\item \mae{第3章 問題2・10-18，10-21，問題5・10-18} 複素数の係数を実数として扱った。$\phi^{*}$ を書く行で $r\to r^{*}$，$t\to t^{*}$ と係数すべてに ${}^{*}$ を付け，2乗は $\abs{r}^{2}$ と書く。
\item 古典との比較は，まず古典の粒子がどうなるかを決める。$E>V_{1}$ の段差は古典では必ず越える（$R=0$，$T=1$）。
\item 有効数字は次の桁で四捨五入し，$R+T=1$ で検算する。
